\begin{textsample}{POS Dim 1 – vem\_ed\_21 – Score 3.00 – vem\_ed\_21/t104.txt}  \label{ex:f1_pos_171}
Debate internacional sobre Intercâmbios Virtuais A conferência IVEC 2023 (International Virtual Exchange Conference) \textbf{é} \textbf{reconhecida} como o mais importante evento sobre Intercâmbios Virtuais no mundo.
Pela primeira vez, ocorreu no Hemisfério Sul, no Brasil, em São Paulo, de 30 de outubro a 1 de novembro de 2023.
A IVEC 2023 reuniu 496 participantes (presenciais e on-line) de 41 países e \textbf{foi} organizada pela Associação Brasileira de Educação Internacional (Faubai) e pela Unesp: https://iveconference.org/.
Ana Cristina Biondo Salomão, coordenadora \textbf{geral} do Programa Brazilian Virtual Exchange (BraVE) da Faubai e assessora de relações internacionais na Unesp, deu as boas-vindas aos conferencistas na cerimônia de abertura.
A plenária de abertura \textbf{foi} proferida pelo keynote speaker Virgílio Almeida (UFMG / Brasil).
Intitulada “Exploring Emerging Technologies Globally: an opportunity \textbf{for} inclusive and collaborative learning” (Explorando tecnologias emergentes globalmente: uma oportunidade para o \textbf{aprendizado} inclusivo e \textbf{colaborativo}), a palestra abordou como os algoritmos estão moldando espaços de colaboração e como preparar os alunos para trabalhar com a inteligência \textbf{artificial}.

% matched lemmas: aprendizado, artificial, colaborativo, geral, reconhecer, ser
\end{textsample}
